\subsection{Capability-Specific Reinforcement Learning}
\label{sec:training}

Given the set of micro-environments $\{\mathcal{E}_c\}_{c \in \mathcal{C}^{*}}$
produced by the environment synthesis pipeline (\S\ref{sec:pipeline}), we train
a separate low-rank adapter~\citep{hu2022lora} $\Delta_c$ for each capability
$c \in \mathcal{C}^{*}$ while keeping the base model $\pi_\theta$ frozen.  This
modular design isolates capability acquisition: each adapter specializes in a
single skill without interfering with others.

\paragraph{Group Relative Policy Optimization.}
We train each adapter $\Delta_c$ using Group Relative Policy Optimization
(GRPO;~\citealt{shao2024deepseekmath}), an on-policy reinforcement learning
algorithm that eliminates the need for a learned value function.  For each
training iteration, the current policy $\pi_{\theta + \Delta_c}$ generates
$G$ \emph{groups} of rollouts.  Within each group $g$, $K$ independent
trajectories $\{\tau_{g,1}, \dots, \tau_{g,K}\}$ are sampled from the
\emph{same} environment seed $s_g$, so all $K$ rollouts start from an
identical initial state and differ only through temperature sampling.  This
shared-seed structure enables a group-relative baseline: for each trajectory
$\tau_{g,k}$ with terminal reward $r_{g,k}$, the advantage is computed as
\begin{equation}
    \hat{A}_{g,k} = \frac{r_{g,k} - \bar{r}_g}{\sigma_g + \epsilon},
    \qquad
    \bar{r}_g = \frac{1}{K} \sum_{j=1}^{K} r_{g,j},
    \qquad
    \sigma_g = \sqrt{\frac{1}{K}\sum_{j=1}^K (r_{g,j} - \bar{r}_g)^2},
\end{equation}
where the normalization by $\sigma_g$ makes the signal scale-invariant across
environments with different reward magnitudes.  All action tokens within a
trajectory receive the same advantage (since all tokens in a turn contribute to
the same terminal outcome).  Groups where all $K$ rollouts receive
identical rewards are discarded, as they produce zero advantage and contribute
no gradient signal.

The policy is updated via a clipped surrogate objective
\begin{equation}
    \mathcal{L}_{\text{GRPO}} = -\frac{1}{|\mathcal{B}|}\sum_{(g,k) \in \mathcal{B}}
    \frac{1}{T_{g,k}} \sum_{t=1}^{T_{g,k}}
    \min\!\bigl(
        \rho_t \, \hat{A}_{g,k},\;
        \text{clip}(\rho_t, 1{-}\epsilon, 1{+}\epsilon)\,\hat{A}_{g,k}
    \bigr),
    \label{eq:grpo}
\end{equation}
where $\rho_t = \pi_{\theta+\Delta_c}(y_t \mid x, y_{<t}) \,/\,
\pi_{\theta+\Delta_c}^{\text{old}}(y_t \mid x, y_{<t})$ is the per-token
importance sampling ratio between the current and rollout-time policies,
$T_{g,k}$ is the number of action tokens, $\mathcal{B}$ is the mini-batch, and
$\epsilon$ is the clipping threshold.  The ratio $\rho_t$ prevents large policy
updates: when $\hat{A}_{g,k} > 0$ and $\rho_t > 1 + \epsilon$, the gradient
is clipped.

\paragraph{On-Policy Rollout Collection.}
At each iteration, we synchronize the current LoRA adapter to a separate
inference server and generate rollouts on-policy.  The current policy interacts
with the micro-environment's procedural task generator: it receives observations,
calls tools, and receives tool results until the episode terminates.  The
programmatic reward function $R_c$ then scores the completed trajectory.
Because rollouts are collected on-policy (from the current $\pi_{\theta +
\Delta_c}$), the importance sampling ratio $\rho_t$ stays close to 1 at the
start of each iteration, ensuring stable gradient estimates.

\paragraph{Training Details.}
We parameterize each adapter as LoRA with rank $r{=}16$ applied to all
attention and feed-forward projections (\texttt{q}, \texttt{k}, \texttt{v},
\texttt{o}, \texttt{gate}, \texttt{up}, \texttt{down}).  We use AdamW with
learning rate $10^{-5}$, group size $K{=}16$, and $G{=}64$ groups per
iteration.  To reduce wasted compute on non-decision-making turns (e.g.,
read-only API calls that gather information), we filter out training samples
whose only action is an information-retrieval tool call, retaining at
least one sample per episode to preserve the reward signal.  Tool results
from earlier turns are truncated to reduce prompt length.  Sequences are
sorted by length within mini-batches to minimize padding overhead.

\paragraph{Multi-Environment Curriculum.}
When training across multiple capabilities simultaneously, we compose
micro-environments into a curriculum by sampling from a weighted mixture:
at each iteration, each group is assigned to capability $c$ with
probability $w_c$, where
$\{w_c\}_{c \in \mathcal{C}^{*}}$ are tunable mixing weights.
This allows concentrating training compute on the capabilities with the
largest performance gaps while maintaining exposure to all skills.

\subsection{Cross-Skill Distillation}
\label{sec:distillation}

The per-capability training stage produces a set of specialist adapters
$\{\Delta_c\}_{c \in \mathcal{C}^{*}}$, each of which excels at its target
capability but lacks the others.  To consolidate all capabilities into a
single model, we perform on-policy cross-skill distillation
(\citealt{glm2025chatglm}; \citealt{skyrl2025}).

For each capability $c$, we treat the corresponding specialist
$\pi_{\theta + \Delta_c}$ as a \emph{teacher}.  A single \emph{student}
adapter $\Delta_{\text{unified}}$ is trained on-policy: the student generates
trajectories from the mixed environment curriculum, and each trajectory is
scored by the teacher specialist corresponding to that trajectory's capability.
Specifically, the teacher provides per-token logprobs
$\log \pi_{\theta+\Delta_c}(y_t \mid x, y_{<t})$
on the student's own generated tokens, yielding a dense per-token training
signal.

We use a PPO-style clipped surrogate objective where the per-token advantage
is derived from the teacher--student logprob gap:
\begin{equation}
    \hat{A}_t^{\text{distill}} =
        \text{sg}\!\bigl[\log \pi_{\text{teacher}}(y_t \mid x, y_{<t})
        - \log \pi_{\text{student}}^{\text{old}}(y_t \mid x, y_{<t})\bigr],
    \label{eq:distill}
\end{equation}
where $\text{sg}[\cdot]$ denotes stop-gradient.  When the teacher assigns
higher probability than the student to a token, the advantage is positive,
encouraging the student to increase that token's probability; when the teacher
assigns lower probability, the advantage is negative.  This per-token signal is
substantially denser than the trajectory-level reward used in GRPO, enabling
effective distillation with a single rollout per seed (group size $K{=}1$)
rather than the $K{=}16$ rollouts required for group-relative advantages.

The student routes each trajectory to the appropriate teacher based on which
micro-environment generated it, so that a structured-data trajectory is scored
by the structured-data specialist, a multi-step trajectory by the multi-step
specialist, and so on.  This per-skill routing enables simultaneous distillation
from all specialists in a single training run.  The resulting unified adapter
$\Delta_{\text{unified}}$ inherits capabilities from all specialists while
maintaining the coherence of a single policy.
